% Owner: root. Working proof independently reviewed by classification_e and boundary_d.
\begin{proposition}[E.2.20 (Proved): identification of the actual analytic cusp]
In the smooth specialization of E.2.10, normalize
the three distinct critical-value sections to
$0,1,\infty$ and write $q=1/t$ near the cusp.
Let $c(s)$ be the holomorphic interior period
lift supplied by E.2.7.  After shrinking the
$(q,s)$-bidisc, its inverse image in the actual
family is biholomorphic, over that bidisc, to
the prescribed marked Mumford cusp family at
parameter $c(s)$.  The isomorphism induces
the given identification of the regular
torus Jacobians and their integral markings.
Its affine origin may be fixed by choosing
a holomorphic section through a smooth
point of the reduced cusp.
\end{proposition}
\begin{proof}
Write $f:\mathcal N\to\Delta_q\times\Delta_s$
for the proper actual cusp neighbourhood and
$\mathcal W=f^{-1}(q=0)$.  E.2.16--E.2.18
give its reduced relative normal-crossing
divisor, its normalized toric hexagon, and
the unchanged opposite-side identifications.
E.2.19 gives an actual holomorphic action
of $H=(\C^*)^2$ over the bidisc.  Its
period marking is the primitive vanishing
lattice
\[
 \Lambda_{\rm tor}
   =\Z\widehat w\oplus\Z\widehat\delta.
 \tag{E.2.20.1}
\]
We reconstruct the neighbourhood from
this action, keeping track of the
extension of its characters.

\emph{The topology and the exact covering.}
The relative coordinates
$q=x,xy,xyz$ give smooth real lifts
of the two real parameter directions
of $s$ which annihilate $dq$ and
are tangent to every cusp stratum.
Away from the cusp such lifts exist
because $f$ is a submersion.
A smooth partition of unity preserves
both these linear requirements.
On the inverse image of a smaller
closed $q$-disc the lifts integrate
along a sufficiently short parameter
path: $q$ remains constant, and
properness prevents escape.
Thus their flows identify the cusp
neighbourhoods, all strata, and
the regular-fibre inclusion maps
by $q$-preserving diffeomorphisms.

For the known Mumford cusp we
check the covering calculation
directly.  Each maximal toric
chart above $|q|<\epsilon$ is
\[
 \{(x,y,z)\in\C^3:|xyz|<\epsilon\},
\]
which is star-shaped and simply
connected.  Every intersection
contains the same connected dense
open torus $(\C^*)^2\times\Delta_q^*$.
Each connected component of an
intersection is open and must meet
that dense set, so the intersection
is connected.  The same holds for
the intersection with a finite
union of preceding charts.
Van Kampen proves simple
connectedness of each finite
union; a loop lies in such a
union by compactness.  The full
toric cover is therefore simply
connected.  The known free,
properly discontinuous period
action has deck group
\[
 \overline\Lambda=\Lambda/\Lambda_{\rm tor}
                 \simeq\Z^2.
\]
On a regular fibre its covering
is $\C^2/\Lambda_{\rm tor}\to
\C^2/\Lambda$.  The regular-fibre
map to the cusp body's fundamental
group is exactly the projection
to $\overline\Lambda$, and the
$q$-meridian dies through the
section $(x,y,z)=(q,1,1)$.
These statements also agree
with the original manuscript's
Proposition~4.7 and Corollary~4.8,
whose original PDF was checked
\cite[p.~30]{AC26Root}; the
preceding argument verifies
the step we use.  The
$q$-preserving transport transfers
this exact description to the
actual neighbourhood.

Take the universal covering
\[
 \widetilde{\mathcal N}\longrightarrow\mathcal N.
\]
Its deck group is
$\overline\Lambda$, and its
central normalized components
form the planar unfolding of
the hexagon.  The $H$ action
lifts.  Indeed the fundamental
loops of $H$ at a central open
orbit lie in one normalized
del Pezzo component, which is
simply connected, and hence die
in $\pi_1(\mathcal N)$.
The identity lift is unique
and satisfies the action law.
Choose a holomorphic section
$p(q,s)$ through a smooth
central cusp point and lift
it to the cover.  On $q\ne0$
its orbit map is an isomorphism
\[
 H\times(\Delta_q^*\times\Delta_s)
       \xrightarrow{\ \sim\ }
 \widetilde{\mathcal N}|_{q\ne0}.
 \tag{E.2.20.2}
\]
The stabilizer has been removed
exactly by (E.2.20.1), using
the just-verified regular-fibre
inclusion map.  In particular
the two characters of $H$
are globally defined on this
open covering space.

Connectedness of $H$ makes
its lifted action commute
with every deck transformation:
conjugation acts continuously
on the discrete deck group
and is the identity at the
identity of $H$.  On
(E.2.20.2), the two deck
generators therefore act by
\[
 (h,q,s)\longmapsto(P_i(q,s)h,q,s),
 \qquad i=1,2.
 \tag{E.2.20.3}
\]
The functions $P_i$ are the
actual marked period multipliers.
Changing the affine origin
does not change them.
The marked period comparison
of E.2.7 and E.2.12 identifies
them with the prescribed
$q$-monomials times holomorphic
units at $c(s)$.

\emph{Character extension across adjacent components.}
For a central component $D_v$,
choose a holomorphic base
section $a_v(q,s)$ through
its dense $H$ orbit.  Its
orbit map is an open
holomorphic embedding
\[
 H\times(\Delta_q\times\Delta_s)
       \longrightarrow\widetilde{\mathcal N}.
\]
It is injective on the
central open orbit and on
each regular fibre, and
its differential is invertible
along those loci.  Its
image contains the whole
regular covering.  Two
such origins differ for
$q\ne0$ by a holomorphic
translation $h_{vw}(q,s)$
in $H$.  We must prove
that this function has no
essential singularity at
$q=0$.

If both origins lie in
the same central component,
their ratio extends as a
holomorphic $H$-valued
unit.  The central open
orbit is free and transitive,
so this follows from the
holomorphic inverse-function
theorem for the orbit map,
with both parameters retained.

For two adjacent components,
use a crossing chart with
$q=ab$ and an orbit coordinate
$z\ne0$.  The stabilizer
of a general conductor
point is a primitive
connected one-dimensional
subtorus of $H$.  Its
two normal characters are
$+1,-1$: on each normalized
toric component these are
the primitive boundary
characters, and invariance
of $q$ makes them opposite.
Compact-circle averaging
gives these exact local
linear coordinates.  A
complementary torus direction
provides the orbit coordinate;
an invariant unit in $q$
can be absorbed while
preserving the weights.

Fix small nonzero constants
$a_0,b_0,z_0$ in this chart
and take the two special
origins
\[
 (a,b,z)=(q/b_0,b_0,z_0),\qquad
 (a,b,z)=(a_0,q/a_0,z_0).
 \tag{E.2.20.4}
\]
The stabilizer cocharacter
evaluated at $a_0b_0/q$
sends the first to the
second.  This equality
is justified by the local
action itself: a path
in that subgroup can
increase the modulus of
$a$ and decrease that
of $b$, with $ab=q$
fixed, staying inside
the crossing chart;
the phase can be varied
without changing either
modulus.  Thus the
possibly large group
parameter causes no
unjustified use of a
local coordinate outside
its domain.

Comparing (E.2.20.4) with
arbitrary chosen origins
adds only holomorphic
$H$-valued units, by the
same-component argument.
The adjacent comparison
therefore has the form
\[
 h_{vw}(q,s)=q^{n_{vw}}u_{vw}(q,s),
 \qquad n_{vw}\in\Z^2,
 \tag{E.2.20.5}
\]
where $u_{vw}$ is holomorphic
with values in $H$ on the
whole bidisc.  Every central
component is reached by
a finite adjacent path.
Multiplying these expressions
proves the same assertion
for its comparison with
the initial origin.
Different paths agree
on the regular cover,
and hence define identical
meromorphic characters.
The exponents are locally
constant in $s$, since
(E.2.20.5) holds with a
nonvanishing unit on the
full bidisc.

\emph{The fan and the analytic coordinate maps.}
At either triple-point
section the whole $H$
fixes the point.
Compact-torus averaging
gives three eigen-coordinates
with the constant integer
characters of E.2.19.
Their only primitive
positive invariant relation
is their product.
The reduced normal-crossing
equation therefore writes
$q$ as that product times
a holomorphic invariant
unit.  The unit is a
function of the product
and $s$; absorb it into
one eigen-coordinate to
obtain exactly $q=xyz$.

The orders of the two
global characters along
$D_v$, together with
$\operatorname{ord}_{D_v}q=1$,
give a ray $(v,1)$.
At a triple point the
three incident rays
form a unimodular cone.
Indeed all adjacent
orders and tangent
characters are constant
in $s$ and coincide
with the known ones
for $s\ne0$.
The entire resulting
fan is consequently
the prescribed $A_2$
triangulation, with
its fixed labels.
No additional refinement
is being introduced.

In a triple chart a
global character is
a Laurent monomial in
$x,y,z$ times a function
of $q,s$.  Evaluate it
on a section keeping
two nonzero coordinates
fixed and solving the
third from $q$.
The origin comparisons
above show that the
remaining function is
a power of $q$ times
a holomorphic unit.
Thus the dual monomial
coordinates of the
unimodular fan cone
pull back to $x,y,z$
times holomorphic units
in $q,s$.  They extend
holomorphically and
give a local
biholomorphism through
the triple point.
The crossing and
open-orbit charts give
the same conclusion
on the other strata.
The maps agree on
the dense regular
cover and therefore
glue.

We finally check the
common neighbourhood
and the deck action.
We do not take an
infimum of infinitely
many chart radii.
Choose finitely many
charts covering one
compact normalized
central del Pezzo
component and finitely
many overlaps for
the two deck generators.
Shrink the bidisc once
for this finite data.
On the regular cover
the character map
intertwines the actual
deck action with the
reference multipliers
$P_i(q,s)$ in
(E.2.20.3).  Both maps
extend in the charts
just constructed, so
their equality extends
to the central strata.
The reference multipliers
preserve the prescribed
invariant fan.  Deck
translates of the finite
chart collection cover
the entire central
unfolding and use the
same base bidisc.

We have obtained local
biholomorphisms at all
central strata, bijections
on their labelled torus
orbits, and the original
isomorphism on the whole
regular cover.  The
inverse coordinate maps
are the same unimodular
monomial charts with
the inverse units.
Quotienting by the
identical free properly
discontinuous deck
actions gives the
claimed biholomorphism
of actual analytic
cusp neighbourhoods.
\end{proof}

\begin{remark}[Scope and source comparison]
The toric covering and its
period action are the known
Mumford construction.  The
point here is to recover
that construction from
the actual smooth limit,
using the action proved
in E.2.19.  In particular
the absence of an essential
singularity in the character
comparison is proved in
(E.2.20.4)--(E.2.20.5);
it is not an assumed
meromorphic comparison.
The local identification
does not alone identify
the global threefold or
settle arbitrary deformation
chains.  Those require
the finite affine charts
and global collar argument,
and an additional propagation
argument for general bases.
\end{remark}
